% =====================================================================
\appendix

\chapter{Supplementary Proofs and Arguments}\label{app:proofs}

This appendix supplies complete proofs, or fuller arguments, for several results in the
main text. Where a complete proof is already given there, it is not repeated.

\section{Notes on Proposition \ref{prop:overconstrain}}

The proof given in the main text rests on the transversality theorem. Two points are added
here.

\paragraph{Why ``independent constraints'' are necessary.}
If two constraints are not independent (for instance, when the second is a consequence of
the first), their contributions to the codimension do not add,
$\mathrm{codim}\,\mathcal{S} < c$, and the solution set has dimension greater than
$d - c$. The conclusion of Proposition \ref{prop:overconstrain} must therefore be read
under the premise that the constraints are independent and transversal. That premise
usually holds in low-level design, where each constraint originates in a distinct
pre-existing object (binary compatibility, the debugger, exception unwinding, the
system-call interface), and often fails in high-level design, where several
``constraints'' are in fact different formulations of one preference.

\paragraph{Why ``solution-set degeneracy'' does not mean ``unique solution''.}
When $c \ge d$ the solution set is finite or empty, and a finite set may contain more than
one element. Low-level design is nonetheless ``nearly unique'' in practice because, over
and above the formal constraints, an implicit economic constraint is at work: maintaining
two incompatible calling conventions at once carries labour and toolchain costs, and that
cost induces a strict preference order among the admissible options.

\section{Complete Argument for Theorem \ref{th:integration}}

The main text gives the principal argument. The choice of cost function and the
reasonableness of the assumptions are supplied here.

\paragraph{The definition of cost.}
Let the unit interaction cost $c$ be the whole of the work required to handle one
interaction between a particular entity and a particular resource correctly: design,
implementation, testing and long-term maintenance. In the no-substrate case every entity
must do this work independently, and must handle the resource invariants independently as
well.

\paragraph{The magnitude of the consistency obligation.}
For a resource $r$, $n$ independent implementations must agree on the invariants of $r$.
If consistency is not enforced, the deviation of any one implementation breaks the global
invariant. If it is enforced, the enforcement itself requires that all $n$ sites know the
invariant, so the number of enforcement sites is $\Omega(n)$. Neither case improves on the
conclusion ``read-set $\ge k$'' (Theorem \ref{th:centralise}). In other words, the absence
of a substrate is not merely more costly: it forces invariants to be scattered, and
thereby \emph{permanently raises} the read-set lower bound of the system. This is the
second benefit of a unified substrate, and it is independent of amortisation.

\paragraph{Relaxing assumption (b).}
Assumption (b), that every entity participates in all $m$ resources, is the worst case. If
entity $e$ participates in only $|\mathcal{R}(e)|$ resources, the no-substrate cost is
$\sum_e |\mathcal{R}(e)|$. As long as $\mathbb{E}_e[|\mathcal{R}(e)|]$ grows with $m$
(that is, resource use is diffuse and broad), the product-form growth still holds. The
qualitative conclusion of \eqref{eq:threshold} therefore does not depend on the worst-case
assumption.

\section{Supplement to the Item-by-Item Argument for Theorem \ref{th:six}}

The main text argues item by item, showing that the omission of any one function yields
undefined behaviour. Two remarks on the strength of that argument are added here.

\paragraph{Strength of the (F3) scheduling argument.}
The argument relies on the claim that $\to$ must be a function, or that its
non-determinism must be explicitly specified. If the semantics of a system is defined as
the set of \emph{all} interleavings, that is, if non-determinism is taken as part of the
semantics and no selection among it is made, then the selection function may indeed be
absent. In that case, however, the \emph{implementation} of the system must still select
one interleaving in every concrete execution, so the scheduling function remains present
at the implementation level even though it does not appear in the abstract semantics. The
formulation used in this paper, ``either in the system's own code or in its host'', covers
this case.

\paragraph{Strength of the (F6) isolation argument.}
The argument relies on fault reachability. If the execution environment of a system
guarantees the absence of faults (pure functional computation, for instance, or execution
in a restricted environment verified to be fault-free), then (F6) may be absent or may
degenerate to the trivial. This is precisely why some minimal systems can do without the
complete six functions, and it shows that the premises of Theorem \ref{th:six} are not
redundant.

\section{The Read-Set Lower Bound of Theorem \ref{th:centralise}}

The main text proves $|\rho(c)| \ge |\mathrm{S}(\iota)|$. Why that inequality is tight is
explained here.

If $\mathrm{S}(\iota) = \{M_0\}$ and the interface of $M_0$ is sufficient
(Definition \ref{def:neardecomp}), then every change that could break $\iota$ must be
expressed at the interface of $M_0$, so an agent performing the change need only
understand $M_0$ and its call sites. Then $|\rho(c)| = 1 + O(\text{call sites})$, and the
lower bound $1$ is attained (when the change itself occurs inside $M_0$).

If $\mathrm{S}(\iota) = \{M_0, M_1\}$ and both modules bear partial responsibility for
$\iota$, the agent must understand both \emph{and the consistency rule between them},
whose comprehension cost is strictly greater than that of the two modules together. The
read-set therefore grows superlinearly in $|\mathrm{S}(\iota)|$, and the linear relation is
only a lower bound. The bound given in the main text is consequently conservative.

\section{Exhaustiveness of the Three Grounds of Proposition \ref{prop:three-grounds}}

This argument has limits. The three grounds take the following form:

\begin{enumerate}[label=(g\arabic*)]
\item Preferences are generative: there exist an intent $I$ and a situation $s$ such that
      some component of $I$ is fixed only after the bearer of the intent encounters $s$.
\item Specifications are lossy compressions: there exist an intent $I$ and a specification
      $\Phi$ such that $\Phi$ cannot carry all information relevant to $I$.
\item Validation does not reduce to verification: there exist a requirement $R$ and an
      implementation $P$ such that ``$P$ satisfies $\Phi$'' is true while ``$P$ satisfies
      $R$'' is false, and the procedure that decides the latter cannot be reduced to the
      procedure that decides the former.
\end{enumerate}

The three grounds are not independent of one another: (g1) is a source of (g2), and (g3)
is a consequence of (g2). A more accurate formulation may therefore be that \emph{the root
of the intent gap is the composite of (g1) and (g2)}. This paper retains the three
parallel formulations because they correspond to different engineering remedies: (g1) to
interactive refinement and counterexample generation, (g2) to traceability and coverage
analysis, and (g3) to an independent validation process and an allocation of
responsibility.

Whether a fourth ground exists is an open problem (Chapter \ref{ch:objections}, items
(O)).

\section{Tightness of the Bound in Theorem \ref{th:condorcet}}

Equation \eqref{eq:chernoff} is the Hoeffding bound, which is not tight for large $n$. For
$p < 1/2$ the probability of a majority error is in fact given by the central limit
theorem, $\Pr[\text{majority errs}] \approx \Phi\bigl(-\sqrt{n}\,(1/2-p)/\sigma\bigr)$,
whose exponent is $\Theta(n (1/2-p)^{2})$ with a better constant.

This paper chooses the Hoeffding bound over a tighter estimate for three reasons:

\begin{enumerate}[label=(\alph*)]
\item it requires no assumption about the distribution of individual errors, only
      independence, boundedness and a mean;
\item its form is simpler, and it is therefore less easily misread as an exact value;
\item the central conclusion of this paper (Proposition \ref{th:correlated}) does not
      depend on the tightness of the bound: the lower bound $q$ given by a correlated
      error mode is independent of the constant in the upper bound.
\end{enumerate}

\section{A Stronger Form of Proposition \ref{th:correlated}}

The main text gives $\Pr[\text{majority errs}] \ge q$. A finer form is available. Let the
common error mode $C$ and the individual independent errors $I_i$ satisfy: verifier $i$
errs if and only if $C = 1$ or $I_i = 1$. The event that the majority errs is then
\[
\{C = 1\} \;\cup\; \Bigl\{C = 0,\; \sum_{i=1}^{n} I_i \ge \tfrac{n}{2}\Bigr\},
\]
so that
\[
\Pr[\text{majority errs}] \;=\; q + (1-q)\,\Pr\Bigl[\sum_{i=1}^n I_i \ge \tfrac{n}{2}\Bigr]
\;\ge\; q .
\]
For fixed $q > 0$, moreover,
$\lim_{n \to \infty} \Pr[\text{majority errs}] = q$. This yields a statement stronger than
a lower bound of $q$: \emph{consensus reliability has an asymptotic floor determined by
the correlated error mode; increasing the number of verifiers can only approach that
floor and cannot break through it}. This is the quantitative basis of
Corollary \ref{cor:diversity}.

\section{Notes on Proposition \ref{prop:active}}

The argument in the main text uses the definition of information gain. One clarification
is added here: Proposition \ref{prop:active} does not say that passive reading is without
value. Passive reading supplies the \emph{prior}, in the Bayesian sense, whereas active
probing supplies the \emph{likelihood}. A prober with no prior does not know what to
probe; a reader with only a prior cannot update. The first four of the seven steps of the
protocol in \S\ref{sec:protocol} are therefore reading (establishing the prior), and the
last three are probing (acquiring the gradient). Their order is not arbitrary.
